% Generated draft: requires mathematical and source review.
\documentclass[11pt]{article}
\usepackage[margin=25mm]{geometry}
\usepackage{fontspec}
\setmainfont{Latin Modern Roman}
\usepackage{amsmath,amssymb,mathtools,unicode-math}
\setmathfont{Latin Modern Math}
\usepackage{enumitem,needspace,xurl,hyperref}
\setlength{\parindent}{0pt}
\setlength{\parskip}{5pt}
\setlength{\emergencystretch}{3em}
\newcommand{\N}{\mathbb{N}}
\newcommand{\Z}{\mathbb{Z}}
\newcommand{\Q}{\mathbb{Q}}
\newcommand{\R}{\mathbb{R}}
\newcommand{\C}{\mathbb{C}}
\newcommand{\F}{\mathbb{F}}
\newcommand{\A}{\mathbb{A}}
\newcommand{\PP}{\mathbb{P}}
\newcommand{\E}{\mathbb{E}}
\newcommand{\eps}{\varepsilon}
\newcommand{\dd}{\,\mathrm{d}}
\begin{document}
\section*{UnsolvedMath ID: 30002311}
\textbf{Source problem code:} catalogue entry
\subsection{Definitions and mathematical statement}
Let \(\varepsilon>0\) denote the microscopic scale and \(h>0\) the coarse discretization scale. The supplied record denotes by \(u_\varepsilon^h\) a multiscale numerical approximation of a stochastic heterogeneous partial differential equation and by \(u_0^h\) its corresponding numerical homogenized solution. Write
\[
F_{\varepsilon,h}=a_{\varepsilon,h}\bigl(u_\varepsilon^h-u_0^h\bigr),
\]
where \(a_{\varepsilon,h}>0\) is a normalization to be specified from the underlying model; no normalization rate is supplied.

The requested problem is to determine, for each fixed admissible coarse scale \(h\), the limiting probability law \(\mu_h\) of \(F_{\varepsilon,h}\) as \(\varepsilon\downarrow0\), and then determine whether these intermediate laws converge as \(h\downarrow0\) to the fluctuation law \(\mu_{\mathrm{corr}}\) predicted by continuum corrector theory. Schematically, the first target is
\[
\mathcal L(F_{\varepsilon,h})\Longrightarrow\mu_h
\qquad(\varepsilon\downarrow0,\ h\text{ fixed}),
\]
where \(\mathcal L\) denotes probability law and the state space and meaning of convergence must be specified. If the discrete fluctuations take values in spaces \(E_h\), comparison with a continuum law requires specified measurable reconstruction maps \(I_h:E_h\to E\) into a common topological space \(E\). With those choices, the second target can be expressed as
\[
(I_h)_\#\mu_h\Longrightarrow\mu_{\mathrm{corr}}
\qquad(h\downarrow0),
\]
where \((I_h)_\#\mu_h\) is the distribution obtained by applying \(I_h\) to a random element with law \(\mu_h\). These spaces and maps are placeholders for missing model data, not additional hypotheses supplied by the source.

The record does not specify the equation, spatial dimension, domain, boundary conditions, coefficient law and correlation assumptions, numerical scheme, any additional discretization parameters, normalization, convergence topology, or precise continuum corrector law. Consequently, these displays describe the requested successive limits but do not constitute a uniquely specified model-level conjecture. Existence of either limit and identification of its law are questions, not assertions. The supplied dated assessment further qualifies the intermediate limit as dependent on the discretization and correlation regime; it does not establish a universal intermediate distribution.
\subsection{Short English statement}
For each fixed coarse mesh scale, identify the limiting law of the suitably rescaled difference between the stochastic numerical solution and its numerical homogenized counterpart. Then ask whether these laws approach the continuum corrector fluctuation law as the mesh is refined. The source leaves the model, scaling, and mode of convergence unspecified, so a precise model-level statement requires further information.
\subsection{Sources}
\begin{itemize}
\item Local registry: the supplied UnsolvedMath problem record, permanent ID \(30002311\), source code \texttt{catalogue entry}; registry files \texttt{data/unsolvedmath/problems.json} and \texttt{catalog/problems.json}. The supplied collection is \emph{Oberwolfach Reports Open Problems}; the proposed year is 2013, and no individual proposer is identified.
\item Supplied report citation: \emph{Interplay of Theory and Numerics for Deterministic and Stochastic Homogenization} (2013), Oberwolfach Reports, DOI 10.4171/owr/2013/14, pp. 801--865. \url{https://doi.org/10.4171/owr/2013/14}. The background describes these reports as records of mathematical workshops and their open-problem discussions. No problem-specific page or speaker attribution is supplied.
\item Supplied literature source: \emph{Oberwolfach Report 14/2013}, \url{https://doi.org/10.4171/owr/2013/14}.
\item Supplied literature source: arXiv:1409.1157, \url{https://arxiv.org/abs/1409.1157}. No title, authors, or supporting passage is supplied in the record.
\item Literature assessment dated exactly \texttt{2026-08-21}: the supplied curation labels the problem partially solved and reports model-dependent intermediate limits, depending on discretization and correlation regime, without establishing a universal intermediate distribution for the full question. This is supplied dated curation derived from the OWR record, not an independently verified assessment; time-sensitive status and cited claims require independent verification.
\item The registry marks the statement \texttt{corrected\_verified} and reports that it was checked against the official Oberwolfach report, preserving the hypotheses, notation, and open target. This is a supplied verification claim; no independent source inspection is asserted here.
\end{itemize}
\end{document}
